\subsection{An indexed minimum budget and a mode-resolution criterion}
\label{ssec:indexed-resolution-budget}

Write $C=\lambda^2$ in this subsection, reserving $c$ for the historical
argument coefficient of the scalar ramp model. The target is one law in
$(C,N,P,k)$ for $1\le k\le1000$, with the first-zero law as its $k=1$
specialization. This is a research target, not a claim of validation over
that complete range. A finite root must first be associated with the
requested ordinal; an unresolved ordinal is an outcome, not a zero-digit
observation.

\paragraph{The indexed minimum.}
For a specified finite branch $\nu_{N,C,k}$, its fixed-cutoff limit
$\nu_{\infty,C,k}$, and the computed root $\widetilde\nu_{N,C,P,k}$, use
the exact signed decomposition
\begin{align}
 e_{\mathrm{num},k}&=(\widetilde\nu_{N,C,P,k}-\nu_{N,C,k})/\gamma_k,\nonumber\\
 e_{\mathrm{mode},k}&=(\nu_{N,C,k}-\nu_{\infty,C,k})/\gamma_k,\nonumber\\
 e_{\mathrm{cut},k}&=(\nu_{\infty,C,k}-\gamma_k)/\gamma_k.
 \label{eq:indexed-signed-channels}
\end{align}
Let $B_{j,k}$ be upper bounds on the absolute values of these three
channels and $D_{j,k}=-\log_{10}B_{j,k}$. Then
\begin{equation}\label{eq:indexed-min-budget}
 D_k\ge-\log_{10}\!\left(\sum_j10^{-D_{j,k}}\right)
 \ge\min\{D_{\mathrm{num},k}(C,N,P),D_{\mathrm{mode},k}(C,N),
                   D_{\mathrm{cut},k}(C)\}-\log_{10}3.
\end{equation}
The first expression is the exact logarithm of the summed upper bounds,
not a fitted smoothing of the minimum. Equality to a minimum up to a
uniformly bounded offset requires two-sided channel estimates and a
quantitative noncancellation condition. Numerical assembly error belongs
in $B_{\mathrm{num},k}$ unless it is separately bounded and counted.

Insufficient $N$ enters through the \emph{indexed mode budget} and through
branch validity. A fourth ``resolution budget'' must not count the same
mode error twice. The finite-parent identity in
Section~\ref{ssec:thousand-zero-campaign} separates eigenvalue, indexed
transfer and remainder increments. If its limiting branches and uniform
bounds are established, their sum supplies a derived inner minimum for
the mode budget. Those estimates have not been proved here.

\paragraph{An exact dual representation of the mode channel.}
The following specialization of goal-oriented Galerkin duality
\cite{IndexedBeckerRannacher} gives a precise object to estimate. At fixed
$C$, let the real self-adjoint operator $A_C$ have a normalized simple
isolated ground state $u$, eigenvalue $\varepsilon_\infty$, and next
spectral point $\lambda_2>\varepsilon_\infty$. Let $P_N$ be the orthogonal
Fourier projection and $u_N$ an exact unit Ritz vector, in the operator
domain, with eigenvalue $\varepsilon_N<\lambda_2$. Choose its sign so
$s_N=\langle u,u_N\rangle>0$. Define
\begin{align*}
 r_N&=(A_C-\varepsilon_N)u_N, &P_Nr_N&=0,\\
 Q&=I-|u\rangle\langle u|,
 &R_N&=Q(A_C|_{u^\perp}-\varepsilon_N)^{-1}Q.
\end{align*}
Here $\|R_N\|\le(\lambda_2-\varepsilon_N)^{-1}$. For real $t$, let
$b(t)$ be the vector representing finite-interval Fourier-transform
evaluation, $F_a(t)=\langle b(t),a\rangle$, and set $z_N(t)=R_Nb(t)$.
For complex $t$ use the analytic coefficient extension and the bilinear
transpose pairing. The apparent lattice poles in the coefficient formula
for $b$ are removable.

\begin{proposition}[Indexed dual mode identity]
\label{prop:indexed-dual-mode}
Under these hypotheses,
\begin{equation}\label{eq:indexed-dual-functional}
 F_{u_N}(t)=s_NF_u(t)
       +\langle (I-P_N)z_N(t),r_N\rangle.
\end{equation}
If $\nu_{N,C,k}$ and $\nu_{\infty,C,k}$ are associated real roots, then
for some real $\xi$ between them,
\begin{equation}\label{eq:indexed-dual-displacement}
 \nu_{N,C,k}-\nu_{\infty,C,k}
 =-\frac{\langle(I-P_N)z_N(\nu_{N,C,k}),r_N\rangle}
              {s_NF_u'(\xi)},
\end{equation}
provided the denominator is nonzero.
\end{proposition}
\begin{proof}
Write $u_N=s_Nu+e$ with $e\perp u$. Projection of the residual onto
$u^\perp$ gives $e=R_Nr_N$. Apply the linear evaluation functional and
use $P_Nr_N=0$ to prove~\eqref{eq:indexed-dual-functional}. At the finite
root this gives $F_u(\nu_{N,C,k})$; the real mean-value theorem between
the associated roots proves~\eqref{eq:indexed-dual-displacement}.
\end{proof}

\begin{proposition}[Sufficient mode-resolution condition]
\label{prop:indexed-resolution-disc}
Let a disc of radius $r$ about a finite simple real root $\nu_{N,C,k}$
contain exactly one zero of $F_{u_N}$, counting multiplicity, and no zero
on its boundary. If
\begin{equation}\label{eq:indexed-resolution-disc}
 \|r_N\|\sup_{|t-\nu_{N,C,k}|=r}\|(I-P_N)z_N(t)\|
  <\min_{|t-\nu_{N,C,k}|=r}|F_{u_N}(t)|,
\end{equation}
then $F_u$ also has exactly one zero in that disc. It is real and simple.
\end{proposition}
\begin{proof}
Apply Rouche's theorem to~\eqref{eq:indexed-dual-functional}; the nonzero
constant $s_N$ does not change zeros. Reality under conjugation forces the
unique zero in a disc centered on the real line to be real.
\end{proof}
If $|F_{u_N}''|\le M_2$ throughout the disc, a usable lower bound for
the right side is
$r|F_{u_N}'(\nu_{N,C,k})|-M_2r^2/2$, when positive. This condition
identifies a limiting-transform root, not its zeta ordinal or the cutoff
error. Disjoint discs and complete root counts are needed for a prefix.

\paragraph{What would justify doubled convergence.}
Suppose $\|r_N\|\le a_N$, $\|(I-P_N)z_N(\nu_{N,C,k})\|\le b_{N,k}$,
$s_N\ge s_*>0$, $|F_u'|\ge m_{N,k}>0$ on the established root interval,
and $0<\Gamma_k\le\gamma_k$. Then
\begin{equation}\label{eq:indexed-dual-budget}
 B_{\mathrm{mode},k}\le\frac{a_Nb_{N,k}}{s_*m_{N,k}\Gamma_k}.
\end{equation}
In particular, conditional envelopes
\[
 a_N\le A(C)e^{-a(C)\Phi(N)},\qquad
 \frac{b_{N,k}}{s_*m_{N,k}\Gamma_k}
       \le B(C,k)e^{-b(C,k)\Phi(N)}
\]
give
\begin{equation}\label{eq:indexed-dual-depth}
 D_{\mathrm{mode},k}\ge
 \frac{[a(C)+b(C,k)]\Phi(N)-\log[A(C)B(C,k)]}{\log10}.
\end{equation}
The exponent doubles only if the dual factor supplies the same decay rate
with controlled constants and domain. Those CCM rates remain open. This
identity does not promote a finite-ramp stretched exponent to an
infinite-$N$ coefficient asymptotic.

\paragraph{Minimum required basis size.}
For requested depth $D_*$, allocate each channel at most
$10^{-D_*}/3$. A sufficient mode threshold is the first $N$ in a proved
envelope domain for which the branch condition above holds and
\begin{equation}\label{eq:indexed-N-res}
 [a(C)+b(C,k)]\Phi(N)
 \ge \log[A(C)B(C,k)]+(D_*+\log_{10}3)\log10.
\end{equation}
This defines a conditional $N_{\mathrm{res}}(C,k,D_*)$; the required
functions are not yet an evaluated closed law. A monotone envelope would
make it a threshold for all larger $N$, whereas an isolated measured
success can result from signed cancellation.

The count, band and accuracy constraints are different. A consecutive
prefix of $K$ movable positive nonlattice roots requires $K\le N$ by the
degree-$N$ secular representation. Thus a movable first-1000 prefix
requires $N\ge1000$. Geometric inclusion below
$2\pi N/\log C$ neither guarantees accuracy nor is necessary for an
isolated accurate match: retained data include successful matches above
that band. Neither count nor band placement supplies
\eqref{eq:indexed-N-res}.

\paragraph{Precision also depends on the selected root.}
Writing $r_u(t)=u\cdot v_N(t)$, where
$v_0(t)=1/t$ and $v_j(t)=\sqrt2\,t/(t^2-\alpha_j^2)$, implicit
differentiation at a simple nonlattice root gives
\begin{equation}\label{eq:indexed-root-condition}
 \mathrm d\nu[u](e)=-\frac{e\cdot v_N(\nu)}{r_u'(\nu)},\qquad
 \chi_N(\nu)=\frac{\|v_N(\nu)\|}{|\nu r_u'(\nu)|},\qquad
 \|v_N(\nu)\|^2=\frac1{\nu^2}
       +2\nu^2\sum_{j=1}^N\frac1{(\nu^2-\alpha_j^2)^2}.
\end{equation}
Because $u\cdot v_N(\nu)=0$, the worst direction is tangent to the
unit-state sphere. If a state error were bounded by $C_0 10^{-P}$,
the local precision depth would be approximately
$P-\log_{10}C_0-\log_{10}\chi_N$; a finite error bound needs interval
derivative and evaluation-vector bounds. Requested digits and matrix bits
do not automatically provide that state-error premise. The point
sensitivities in the tested high-index stencils have
$\log_{10}\chi_N\simeq672.6$--$676.2$ at $(C,N)=(1000,890)$ and
$734.0$--$745.8$ at $(1200,970)$. They describe arbitrary coefficient
perturbations, not measured corruption or the structured mode error.

\paragraph{Status.}
The identities and implications above have explicit spectral and branch
hypotheses. Twenty-eight exact rational self-adjoint fixture checks verify
their finite algebra. Uniform arithmetic residual, dual-tail, derivative
and cutoff bounds, with validation across independently qualified source families,
remain proof and data targets. Section~\ref{ssec:thousand-zero-campaign}
now supplies a complete counted first-1000 test on one stored state.
